\documentclass[11pt]{article}

% ============================================================
% GUÍA-RÚBRICA · EVALUACIÓN DE PARES DE EXPOSICIÓN
% ISR-401 Ingeniería de Requerimientos · SPA 2026-2027
% Aplica a: GA-SUM-06 (EXP-2, S13) — Herramientas CASE e Estándares de IR
% Compilación: pdflatex -> pdflatex
% ============================================================

\usepackage[utf8]{inputenc}
\usepackage[T1]{fontenc}
\usepackage[scaled]{helvet}
\renewcommand{\familydefault}{\sfdefault}
\usepackage[a4paper,top=2.6cm,bottom=2.3cm,left=2cm,right=2cm,headheight=34pt]{geometry}
\usepackage{graphicx}
\usepackage{array}
\usepackage{tabularx}
\usepackage{longtable}
\usepackage{booktabs}
\usepackage{colortbl}
\usepackage{xcolor}
\usepackage{enumitem}
\usepackage{fancyhdr}
\usepackage{lastpage}
\usepackage[most]{tcolorbox}
\usepackage{amsmath}
\usepackage{caption}
\usepackage[hidelinks]{hyperref}
\usepackage{multirow}

% ---------- Textos en español (sin babel) ----------
\renewcommand{\tablename}{Tabla}
\renewcommand{\refname}{Referencias}

% ---------- Paleta institucional ----------
\definecolor{utegreen}{HTML}{0B7A3B}
\definecolor{utegreendk}{HTML}{075C2C}
\definecolor{midblue}{HTML}{1F5C8B}
\definecolor{gold}{HTML}{D4A017}
\definecolor{rowgray}{HTML}{EFEFEF}
\definecolor{boxgray}{HTML}{F4F6F4}
\definecolor{lightred}{HTML}{FBEAEA}
\definecolor{lightgold}{HTML}{FBF3DC}
\definecolor{lightblue}{HTML}{EAF1F7}
\definecolor{lightgreen}{HTML}{DCEFE3}

% ---------- Columnas personalizadas ----------
\newcolumntype{L}[1]{>{\raggedright\arraybackslash}p{#1}}
\newcolumntype{C}[1]{>{\centering\arraybackslash}p{#1}}
\newcolumntype{R}[1]{>{\raggedleft\arraybackslash}p{#1}}
\newcolumntype{Y}{>{\raggedright\arraybackslash}X}

% ---------- Encabezado y pie ----------
\pagestyle{fancy}
\fancyhf{}
\lhead{\footnotesize UTEQ $\cdot$ Ingeniería de Software $\cdot$ Ingeniería de Requerimientos (ISR-401)}
\rhead{\footnotesize Guía-Rúbrica de Evaluación de Pares · EXP-2 (2.° corte)}
\lfoot{\footnotesize Prof. Gleiston C. Guerrero-Ulloa, PhD}
\rfoot{\footnotesize Página \thepage\ de \pageref{LastPage}}
\renewcommand{\headrulewidth}{0.6pt}
\renewcommand{\footrulewidth}{0.4pt}

\setlength{\parindent}{0pt}
\setlength{\parskip}{5pt}
\renewcommand{\arraystretch}{1.28}

% ---------- Atajos ----------
\newcommand{\thdr}[1]{\cellcolor{utegreen}\textcolor{white}{\textbf{#1}}}
\newcommand{\bhdr}[1]{\cellcolor{midblue}\textcolor{white}{\textbf{#1}}}

\newtcolorbox{cajaclave}[1][]{colback=boxgray,colframe=utegreen,boxrule=0.8pt,arc=2mm,fonttitle=\bfseries,#1}
\newtcolorbox{cajapiso}[1][]{colback=lightred,colframe=red!60!black,boxrule=1pt,arc=2mm,fonttitle=\bfseries,breakable,#1}
\newtcolorbox{cajaaviso}[1][]{colback=lightgold,colframe=gold!80!black,boxrule=0.8pt,arc=2mm,fonttitle=\bfseries,breakable,#1}

\begin{document}

% ============================================================
% PORTADA
% ============================================================
\thispagestyle{empty}
\noindent\includegraphics[width=\linewidth]{logoband.png}

\vspace{20pt}
\begin{center}
	{\large Facultad de Ciencias de la Computación}\\[2pt]
	{\large Carrera de Software}\\[30pt]
	{\LARGE\textbf{\textcolor{utegreendk}{GUÍA-RÚBRICA}}}\\[8pt]
	{\Large\textbf{Evaluación de Pares --- Informe de Exposición (2.° corte)}}\\[6pt]
	{\large Ingeniería de Requerimientos (ISR-401)}\\[2pt]
	{\large Segundo Periodo Académico Ordinario (SPA) 2026--2027}
\end{center}

\vspace{20pt}
\begin{center}
	\begin{tabular}{L{0.30\linewidth}L{0.55\linewidth}}
		\textbf{Aplica a} & GA-SUM-06 (EXP-2, semana 13, 18--24 enero 2027)\\
		\textbf{Tema de la exposición} & Herramientas CASE para IR y estándares\\
		\textbf{Tipo de actividad} & GA-EXP (Gestión de Aula --- Exposición grupal)\\
		\textbf{Rol del estudiante} & Evaluador individual de los equipos que NO son el suyo\\
		\textbf{Docente} & Gleiston Cicerón Guerrero Ulloa, PhD\\
		\textbf{Correo institucional} & \texttt{gguerrero@uteq.edu.ec}\\
	\end{tabular}
\end{center}

\vspace{18pt}
\noindent\textcolor{gold}{\rule{\linewidth}{1.5pt}}
\begin{center}
	\small La \emph{nota de la exposición} la asigna el docente; la nota de este \emph{informe de evaluación de pares} se registra de forma independiente como componente GA de la actividad GA-SUM-06. Ambas contribuyen al cálculo del 2.° corte.
\end{center}

\newpage

% ============================================================
\section{\textcolor{utegreendk}{Propósito y contexto}}
% ============================================================

En la semana 13, cada equipo expone sobre herramientas CASE aplicadas a la IR y sobre los estándares que las rigen (IEEE 29148:2018, ISO/IEC\,25010:2023). Los integrantes de los demás equipos actúan como evaluadores pares: observan, registran evidencias y emiten un juicio técnico fundamentado. Este ejercicio desarrolla la capacidad de evaluar herramientas de software con criterios objetivos y prepara al estudiante para seleccionar herramientas en proyectos reales.

El \textbf{informe de evaluación de pares} es el documento escrito que el evaluador entrega al finalizar la sesión. Cada estudiante entrega \textbf{un informe por cada equipo observado} (excluido el propio).

\begin{cajaclave}[title=Regla de exclusión]
	El estudiante \textbf{no evalúa} a los integrantes de su propio equipo. Si el aula tiene cuatro equipos (A, B, C, D) y el evaluador pertenece al equipo A, entrega informes independientes para los equipos B, C y D.
\end{cajaclave}

\begin{cajaclave}[title=Qué NO debe hacer el evaluador]
	\begin{itemize}[leftmargin=14pt,itemsep=2pt]
		\item No filma ni graba audio o video a los expositores.
		\item No interrumpe la exposición (salvo en el turno de preguntas habilitado por el docente).
		\item No consulta con sus compañeros durante la observación para acordar puntuaciones.
		\item No copia ni parafrasea el informe de otro evaluador.
	\end{itemize}
\end{cajaclave}

% ============================================================
\section{\textcolor{utegreendk}{Contexto temático: qué debe cubrir EXP-2}}
% ============================================================

El equipo expositor puede organizar su presentación en torno a uno o más subtemas de la Unidad IV. La siguiente lista orienta al evaluador par sobre qué elementos \emph{pueden} y \emph{deben} aparecer en una exposición de alta calidad sobre este tema:

\begin{itemize}[leftmargin=16pt,itemsep=3pt]
	\item \textbf{Herramientas CASE de gestión ágil:} Jira (epics, historias, backlogs, tableros Kanban/Scrum) y Trello (tableros, listas, tarjetas con criterios de aceptación).
	\item \textbf{Herramientas CASE de modelado:} Enterprise Architect / Sparx EA (diagramas UML de casos de uso, actividad, clases; trazabilidad de requisitos; generación de documentación SRS).
	\item \textbf{Estándar de intercambio de requisitos:} ReqIF (Requirements Interchange Format) --- formato XML para migrar requisitos entre herramientas.
	\item \textbf{Estándares normativos:} IEEE 29148:2018 (proceso de IR, atributos de un buen requisito, estructura SyRS/SRS) e ISO/IEC\,25010:2023 (modelo de calidad del producto software; características y subcaracterísticas relevantes para la especificación de requisitos de calidad).
	\item \textbf{Técnicas de validación y gestión:} inspección de Fagan, walkthroughs, CCB (Control Board de cambios), trazabilidad, baselines.
	\item \textbf{IR ágil:} product backlog, BDD (Behavior-Driven Development), historias de usuario en contexto SCRUM o XP.
\end{itemize}

\begin{cajaaviso}[title=Para el evaluador par]
	No todos los equipos cubrirán todos los puntos anteriores. Su tarea es evaluar la \emph{calidad y profundidad} de lo que el equipo sí presentó, no penalizar por no haber cubierto subtemas que el equipo no declaró en su alcance. La comparación entre al menos dos herramientas o estándares sí es un indicador esperado en cualquier enfoque.
\end{cajaaviso}

% ============================================================
\section{\textcolor{utegreendk}{Estructura del informe de evaluación de pares}}
% ============================================================

El informe se entrega en papel o en PDF al docente al finalizar la sesión, según se indique. Extensión: \textbf{una a dos páginas por equipo evaluado}. Debe contener:

\subsection{\textcolor{midblue}{Sección 1 --- Identificación}}

\noindent\begin{tabularx}{\linewidth}{|L{0.38\linewidth}|Y|}
	\hline
	\bhdr{Campo} & \bhdr{Descripción}\\
	\hline
	Nombre completo del evaluador & Quién escribe el informe.\\
	\hline
	\rowcolor{rowgray} Equipo del evaluador & Nombre o número del equipo al que pertenece.\\
	\hline
	Equipo evaluado & Nombre o número del equipo que expuso.\\
	\hline
	\rowcolor{rowgray} Herramienta o subtema principal expuesto & Herramienta(s) o norma(s) que el equipo declaró como eje central.\\
	\hline
	Actividad & GA-SUM-06 (EXP-2).\\
	\hline
	\rowcolor{rowgray} Fecha & Día en que se realizó la exposición.\\
	\hline
\end{tabularx}

\subsection{\textcolor{midblue}{Sección 2 --- Registro de indicadores observados}}

Para cada criterio de la Tabla~\ref{tab:criterios}, el evaluador registra:
\begin{enumerate}[leftmargin=18pt,itemsep=3pt]
	\item \textbf{Nivel asignado} (0, 1, 2, 3 o 4) según los descriptores de la Tabla~\ref{tab:niveles}.
	\item \textbf{Evidencia concreta:} al menos una oración con un observable medible (número de diapositiva, nombre exacto de la herramienta o cláusula de norma, dato comparado, pregunta técnica planteada, error conceptual detectado). Frases como «fue bueno» o «estuvo bien explicado» no son evidencias válidas.
\end{enumerate}

\subsection{\textcolor{midblue}{Sección 3 --- Nota sugerida calculada}}

El evaluador aplica la fórmula de la sección~\ref{sec:formula} y anota el resultado.

\subsection{\textcolor{midblue}{Sección 4 --- Retroalimentación constructiva}}

Dos o tres frases dirigidas al equipo expositor: qué fue lo más valioso (técnicamente) y qué aspecto concreto mejoraría más su demostración o análisis de herramientas. Sin adjetivos vacíos; cada frase identifica un elemento específico y medible.

% ============================================================
\section{\textcolor{utegreendk}{Criterios de evaluación de la exposición (EXP-2)}}
% ============================================================

Los criterios de la Tabla~\ref{tab:criterios} describen lo que el evaluador par observa. Los pesos siguen la distribución GA-EXP del encuadre: D1 = 10, D2 = 20, D3 = 15, D4 = 20, D5 = 25, D6 = 10.

\begin{cajaaviso}[title=Indicadores observables y medibles]
	Cada descriptor incorpora indicadores específicos del tema EXP-2 (herramientas CASE e estándares). Al asignar el nivel, cite en la evidencia al menos un observable concreto: nombre de la herramienta, cláusula del estándar, número de criterios de comparación, diapositiva, respuesta en preguntas. Si un indicador no se corresponde exactamente con ningún nivel, aplique el más cercano y justifíquelo.
\end{cajaaviso}

{\footnotesize\setlength{\tabcolsep}{3pt}
\begin{longtable}{|C{0.45cm}|L{2.4cm}|C{0.65cm}|L{3.1cm}|L{3.1cm}|L{3.1cm}|L{2.4cm}|}
	\caption{Criterios de evaluación de la exposición EXP-2 para el evaluador par.}\label{tab:criterios}\\
	\hline
	\thdr{\#} & \thdr{Criterio} & \thdr{Peso} & \thdr{4 --- Excelente} & \thdr{3 --- Bueno} & \thdr{2/1 --- Aceptable/Insuficiente} & \thdr{0 --- Ausente}\\
	\hline
	\endfirsthead
	\hline
	\thdr{\#} & \thdr{Criterio} & \thdr{Peso} & \thdr{4 --- Excelente} & \thdr{3 --- Bueno} & \thdr{2/1 --- Aceptable/Insuficiente} & \thdr{0 --- Ausente}\\
	\hline
	\endhead
	\hline
	\multicolumn{7}{r}{\footnotesize\textit{Continúa en la página siguiente\ldots}}\\
	\endfoot
	\hline
	\endlastfoot
	D1 & \textbf{Cumplimiento de la guía} (alcance y estructura declarados) & 10 & Cubre el alcance declarado al inicio: herramienta(s) presentada(s), comparación con al menos una alternativa, referencia a un estándar aplicable, y demostración o ejemplo concreto, todo presente. & Falta un elemento menor del alcance (p.\,ej., la comparación es superficial o el estándar solo se menciona de paso). & \textbf{(2)} Faltan dos elementos del alcance o la demostración está ausente. \textbf{(1)} Solo expone características de una sola herramienta sin comparación ni estándar. & No expone o el tema no corresponde a herramientas CASE ni estándares de IR.\\
	\hline
	\rowcolor{rowgray} D2 & \textbf{Fundamentación conceptual} (precisión técnica y dominio de herramientas y normas) & 20 & Nombra correctamente $\geq$3 conceptos técnicos (p.\,ej., cláusulas de IEEE 29148, atributos de calidad de ISO/IEC 25010, artefactos ReqIF, funciones de EA); las afirmaciones son verificables y coherentes con los estándares citados. & Conceptos mayoritariamente correctos; $\leq$1 imprecisión que no invalida el contenido central. & \textbf{(2)} Dos o más imprecisiones técnicas o confusión entre funciones de herramientas distintas. \textbf{(1)} Errores conceptuales graves o uso incorrecto sistemático de términos del estándar. & No hay fundamentación técnica o el contenido no pertenece a IR.\\
	\hline
	D3 & \textbf{Desarrollo técnico} (demostración o ilustración de la herramienta en contexto de IR) & 15 & Muestra al menos una función concreta de la herramienta aplicada a un caso de IR (p.\,ej., trazabilidad en EA, tablero Scrum en Jira con criterios de aceptación, diagrama de requisitos en Sparx, archivo ReqIF comentado); los datos o capturas son identificables y tienen fuente. & Muestra la herramienta con un ejemplo básico; el contexto de IR es reconocible aunque poco elaborado. & \textbf{(2)} Describe la herramienta en abstracto sin mostrar ni ilustrar su uso. \textbf{(1)} Solo enumera menús o características sin aplicarlas a ningún artefacto de IR. & No hay demostración ni ilustración técnica evaluable.\\
	\hline
	\rowcolor{rowgray} D4 & \textbf{Análisis crítico y recomendación} (comparación fundamentada y posicionamiento) & 20 & Compara $\geq$2 herramientas o estándares con $\geq$3 criterios explícitos (p.\,ej., soporte a trazabilidad, curva de aprendizaje, integración con otros artefactos, costo/licencia, conformidad con IEEE 29148); la recomendación final identifica el tipo de proyecto para el que cada opción es más adecuada. & Comparación con $\geq$2 criterios; recomendación presente aunque sin especificar tipo de proyecto o contexto. & \textbf{(2)} Comparación descriptiva sin criterios explícitos o recomendación genérica no argumentada. \textbf{(1)} No hay comparación o la recomendación contradice el análisis previo. & Ausente.\\
	\hline
	D5 & \textbf{Comunicación oral y visual} (claridad, tiempo, manejo de preguntas técnicas) & 25 & Exposición fluida sin leer pantalla; tiempo dentro del rango $\pm$2 min; responde $\geq$1 pregunta técnica sobre la herramienta o el estándar con precisión; diapositivas con capturas legibles y textos de fuente $\geq$18 pt visibles desde el fondo. & Uno de los cuatro indicadores cumplido parcialmente (pequeña lectura, tiempo $\pm$3 min, respuesta con imprecisión menor, una captura ilegible). & \textbf{(2)} Dos indicadores parciales. \textbf{(1)} Tres o más indicadores parciales o no responde ninguna pregunta técnica. & No expone o la exposición es ininteligible.\\
	\hline
	\rowcolor{rowgray} D6 & \textbf{Fuentes y créditos} (atribución de capturas, datos y estándares) & 10 & Cada captura de herramienta, dato estadístico o cláusula de norma tiene fuente visible en la diapositiva (nombre de la herramienta + versión, norma + año, URL o manual); las fuentes son verificables. & Mayoría de fuentes presentes; $\leq$2 capturas o datos sin atribuir pero el contenido es correcto. & \textbf{(2)} Varias capturas o afirmaciones sin fuente. \textbf{(1)} Solo una o ninguna fuente visible en todo el juego de diapositivas. & No hay ninguna fuente o las capturas son de fuente desconocida e inverificable.\\
	\hline
\end{longtable}
}

% ============================================================
\section{\textcolor{utegreendk}{Niveles de desempeño}}\label{sec:niveles}
% ============================================================

\noindent{\small\begin{tabularx}{\linewidth}{|C{0.07\linewidth}|L{0.18\linewidth}|Y|}
	\hline
	\thdr{Nivel} & \thdr{Denominación} & \thdr{Descripción operacional para el evaluador par}\\
	\hline
	4 & Excelente & Cumple íntegramente todos los indicadores del descriptor; no se detecta ninguna omisión ni imprecisión técnica en lo observado.\\
	\hline
	\rowcolor{rowgray} 3 & Bueno & Cumple la mayor parte; hay una deficiencia menor que no afecta la comprensión ni la validez técnica.\\
	\hline
	2 & Aceptable & Cumple aproximadamente la mitad; hay deficiencias que limitan pero no anulan la comprensión técnica.\\
	\hline
	\rowcolor{rowgray} 1 & Insuficiente & Cumple menos de la mitad o hay errores graves que comprometen la comprensión técnica.\\
	\hline
	0 & Ausente & El criterio no se puede evaluar porque el elemento no se presentó o no fue observable.\\
	\hline
\end{tabularx}}\label{tab:niveles}

% ============================================================
\section{\textcolor{utegreendk}{Cálculo de la nota sugerida}}\label{sec:formula}
% ============================================================

El evaluador par calcula una \textbf{nota sugerida} que el docente toma como insumo (no como nota definitiva). La nota definitiva de la exposición la asigna el docente.

\[
	\text{Nota sugerida} = \frac{\displaystyle\sum_{i=1}^{6} w_i \times n_i}{4} \times \frac{10}{100}
\]

donde $w_i$ es el peso de la dimensión $i$ y $n_i$ el nivel asignado (0 a 4). Los pesos son: D1 = 10, D2 = 20, D3 = 15, D4 = 20, D5 = 25, D6 = 10.

\begin{cajaclave}[title=Ejemplo de cálculo]
	Niveles observados: D1 = 3, D2 = 4, D3 = 4, D4 = 3, D5 = 3, D6 = 3.\\[4pt]
	$\text{Nota sugerida} = \dfrac{10 \cdot 3 + 20 \cdot 4 + 15 \cdot 4 + 20 \cdot 3 + 25 \cdot 3 + 10 \cdot 3}{4} \times \dfrac{10}{100} = \dfrac{310}{4} \times 0{,}1 = 7{,}75$
\end{cajaclave}

% ============================================================
\section{\textcolor{utegreendk}{Criterios de calificación del informe de evaluación de pares}}
% ============================================================

El \emph{informe} que entrega el evaluador es en sí mismo una actividad evaluada (GA-SUM-06). Se califica sobre 10 con los criterios de la Tabla~\ref{tab:rubrica_informe}.

\begin{cajapiso}[title={Criterios de piso del informe: incumplimiento implica CERO en el informe}]
	\begin{enumerate}[label=\textbf{P\arabic*.},leftmargin=24pt,itemsep=3pt]
		\item \textbf{Entrega al docente al finalizar la sesión.} No se acepta entrega posterior, salvo causa debidamente justificada antes del inicio de la actividad.
		\item \textbf{Un informe por cada equipo observado} (excluido el propio). La ausencia del informe de cualquiera de los equipos observados implica CERO en toda la actividad.
		\item \textbf{Evidencia concreta por criterio.} Si más de dos criterios carecen de evidencia observable (únicamente nivel numérico sin texto), la nota del informe es CERO.
		\item \textbf{Sin copia de otro evaluador.} Informes con evidencias idénticas o notoriamente similares (misma redacción, mismas capturas descritas) reciben CERO en ambos evaluadores implicados.
	\end{enumerate}
\end{cajapiso}

{\footnotesize\setlength{\tabcolsep}{3pt}
\begin{longtable}{|L{2.6cm}|C{0.85cm}|L{3.3cm}|L{3.0cm}|L{3.3cm}|L{2.2cm}|}
	\caption{Rúbrica de calificación del informe de evaluación de pares (EXP-2).}\label{tab:rubrica_informe}\\
	\hline
	\thdr{Criterio del informe} & \thdr{Peso} & \thdr{4 --- Excelente} & \thdr{3 --- Bueno} & \thdr{2/1 --- Acept./Insuf.} & \thdr{0 --- Ausente}\\
	\hline
	\endfirsthead
	\hline
	\thdr{Criterio del informe} & \thdr{Peso} & \thdr{4 --- Excelente} & \thdr{3 --- Bueno} & \thdr{2/1 --- Acept./Insuf.} & \thdr{0 --- Ausente}\\
	\hline
	\endhead
	\hline
	\multicolumn{6}{r}{\footnotesize\textit{Continúa\ldots}}\\
	\endfoot
	\hline
	\endlastfoot
	\textbf{I1. Completitud} (todos los campos y criterios completos) & 20 & Identificación completa (incluye herramienta/subtema expuesto); los 6 criterios con nivel y evidencia; nota sugerida calculada; retroalimentación presente. & Falta un campo menor de identificación o un criterio sin evidencia. & \textbf{(2)} Dos criterios sin evidencia o falta la sección de retroalimentación. \textbf{(1)} Tres o más criterios sin evidencia. & No entregó o entregó una hoja en blanco.\\
	\hline
	\rowcolor{rowgray} \textbf{I2. Objetividad y base observable} (evidencias concretas referenciadas al tema de herramientas CASE) & 30 & Cada nivel está respaldado por al menos una observación concreta y técnicamente específica (nombre de la herramienta, cláusula del estándar, número de diapositiva, función demostrada, pregunta respondida o no respondida); sin afirmaciones vacías como «fue bueno» o «explicaron bien». & Al menos 4 de 6 criterios tienen evidencia concreta técnica; los otros son superficiales. & \textbf{(2)} Solo 2--3 criterios con evidencia concreta. \textbf{(1)} Evidencias todas genéricas u opinativas. & Sin ninguna evidencia observable.\\
	\hline
	\textbf{I3. Coherencia entre nivel y evidencia} (el nivel asignado se justifica con lo descrito) & 25 & En todos los criterios el nivel asignado es coherente con la evidencia descrita (no asigna 4 cuando describe una demostración incompleta, ni 1 cuando describe un ejemplo funcional). & En $\leq$1 criterio el nivel y la evidencia no son del todo coherentes. & \textbf{(2)} En 2--3 criterios hay incoherencia nivel--evidencia. \textbf{(1)} En 4 o más criterios el nivel no concuerda con lo descrito. & Sin coherencia verificable.\\
	\hline
	\rowcolor{rowgray} \textbf{I4. Cálculo correcto} (nota sugerida calculada correctamente) & 10 & La nota sugerida se calcula con la fórmula correcta, los pesos D1--D6 correctos y sin error aritmético. & Error aritmético menor ($\leq$0,2 puntos de diferencia) pero fórmula y pesos correctos. & \textbf{(2)} Error en un peso o en la fórmula. \textbf{(1)} Nota sugerida ausente o calculada sin fórmula. & Sin nota calculada.\\
	\hline
	\textbf{I5. Retroalimentación constructiva} (específica, accionable y técnica) & 15 & Identifica un punto fuerte y un aspecto de mejora concretos referenciados a la herramienta o estándar demostrado (p.\,ej., «la comparación Jira--Trello usó solo precio como criterio; añadir soporte a trazabilidad la haría más técnica»). & Uno de los dos (fortaleza o mejora) es técnicamente concreto; el otro es genérico. & \textbf{(2)} Ambos presentes pero genéricos. \textbf{(1)} Solo uno de los dos presente. & Ausente.\\
	\hline
\end{longtable}
}

Fórmula del informe:
\[
	\text{Nota del informe} = \frac{20 \cdot n_{I1} + 30 \cdot n_{I2} + 25 \cdot n_{I3} + 10 \cdot n_{I4} + 15 \cdot n_{I5}}{4} \times \frac{10}{100}
\]

% ============================================================
\section{\textcolor{utegreendk}{Plantilla de entrega (llenar a mano o en texto)}}
% ============================================================

\begin{tcolorbox}[colback=lightblue,colframe=midblue,boxrule=0.8pt,arc=2mm,title=Informe de Evaluación de Pares --- EXP-2 (GA-SUM-06) --- ISR-401]
	\small
	\noindent\textbf{Evaluador:}~\underline{\hspace{5cm}}\hfill\textbf{Equipo del evaluador:}~\underline{\hspace{2.8cm}}\\[5pt]
	\noindent\textbf{Equipo evaluado:}~\underline{\hspace{4cm}}\hfill\textbf{Actividad:}~GA-SUM-06\hfill\textbf{Fecha:}~\underline{\hspace{2cm}}\\[5pt]
	\noindent\textbf{Herramienta/subtema principal expuesto:}~\underline{\hspace{8.5cm}}\\[8pt]

	\renewcommand{\arraystretch}{1.5}
	\begin{tabularx}{\linewidth}{|C{0.7cm}|L{2.9cm}|C{0.9cm}|Y|}
		\hline
		\textbf{D} & \textbf{Criterio} & \textbf{Nivel} & \textbf{Evidencia concreta observada (herramienta, estándar, dato)}\\
		\hline
		D1 & Cumplimiento guía & & \\[12pt]
		\hline
		D2 & Fundamentación técnica & & \\[12pt]
		\hline
		D3 & Demostración de la herramienta & & \\[12pt]
		\hline
		D4 & Análisis crítico y comparación & & \\[12pt]
		\hline
		D5 & Comunicación oral y visual & & \\[12pt]
		\hline
		D6 & Fuentes y créditos & & \\[12pt]
		\hline
		\multicolumn{2}{|r|}{\textbf{Nota sugerida calculada:}} & \multicolumn{2}{l|}{\hspace{0.5cm}/10}\\
		\hline
	\end{tabularx}

	\vspace{6pt}
	\textbf{Retroalimentación al equipo expositor:}\\[4pt]
	Punto fuerte (técnico): \underline{\hspace{11.5cm}}\\[8pt]
	Aspecto a mejorar (técnico): \underline{\hspace{11cm}}
\end{tcolorbox}

\end{document}
